\documentclass[11pt,a4paper]{article}
\usepackage{auditstyle}

\newcommand{\eps}{\varepsilon}
\newcommand{\W}{\mathcal W}
\newcommand{\dd}{\,d}
\newtheorem{corollary}[theorem]{Corollary}

\title{Part II interacting Weyl reduction\\[3pt]
\large The exact scalar convergence criterion and the excluded quartic LR citation}
\author{Lamport-form proof and proof-obligation ledger}
\date{9 August 2026}

\begin{document}
\maketitle

\begin{resultbox}
\textbf{Reduction completed; interacting convergence remains open.}  To obtain a full
fixed-coarse weak-star scaling limit, it is sufficient to prove convergence of the
single scalar characteristic functional
\[
 \chi_{N,M}(\xi):=\omega_N(\alpha_M^N(W_M(\xi)))
\]
for every fixed \(M\) and every coarse one-particle vector \(\xi\).  Positivity,
boundedness, extension from Weyl polynomials, projective consistency, and continuum
identification are then automatic by the proved steps below.

This static reduction requires neither an interacting Lieb--Robinson bound nor a
uniform spectral gap.  Those inputs are needed only for a quantitative spatial-cutoff
or dynamical theorem.
\end{resultbox}

\tableofcontents

\section{Characteristic-functional criterion}

Let \((\W_N,\alpha_M^N)\) be the Weyl \(C^*\)-inductive system fixed in the
model sheet.  Let
\[
 \mathcal D_M:=\operatorname{span}\{W_M(\xi):\xi\in h_{M,L}\}.   \tag{1.1}
\]
By the definition of the Weyl \(C^*\)-algebra,
\[
 \overline{\mathcal D_M}^{\|\cdot\|}=\W_M.                     \tag{1.2}
\]

\begin{theorem}[full-sequence Weyl criterion]\label{thm:weylcriterion}
For every \(N\), let \(\omega_N\in S(\W_N)\).  Assume that, for every fixed
\(M\) and every \(\xi\in h_{M,L}\), the limit
\[
 \chi_M(\xi):=
 \lim_{N\to\infty,\ N\ge M}
 \omega_N(\alpha_M^N(W_M(\xi)))                               \tag{1.3}
\]
exists.  Then:
\begin{enumerate}
\item there is a unique state \(\rho_M\in S(\W_M)\) satisfying
      \(\rho_M(W_M(\xi))=\chi_M(\xi)\);
\item for every \(A\in\W_M\),
      \[
       \lim_{N\to\infty}\omega_N(\alpha_M^N(A))=\rho_M(A);     \tag{1.4}
      \]
\item the family is projectively consistent,
      \[
       \rho_M=\rho_{M'}\circ\alpha_M^{M'},\qquad M\le M';      \tag{1.5}
      \]
\item there is a unique state \(\Omega_\infty\in S(\W_\infty)\) with
      \[
       \Omega_\infty\circ\alpha_M^\infty=\rho_M.               \tag{1.6}
      \]
\end{enumerate}
\end{theorem}

\section{Proof of Theorem \ref{thm:weylcriterion}}

\LS{1}{1}{Define the limit on finite Weyl polynomials.}

For
\[
 D=\sum_{j=1}^s c_jW_M(\xi_j)\in\mathcal D_M,                  \tag{2.1}
\]
set
\[
 \rho_M^{(0)}(D):=\sum_{j=1}^sc_j\chi_M(\xi_j).                \tag{2.2}
\]
This is independent of the representation (2.1).  Indeed, if \(D=0\) in
\(\W_M\), then \(\alpha_M^N(D)=0\) for every \(N\), and hence
\[
 \sum_jc_j\chi_M(\xi_j)
 =\lim_N\omega_N(\alpha_M^N(D))=0.                             \tag{2.3}
\]
Linearity is immediate.

\LS{1}{2}{Prove boundedness and positivity before extending.}

For \(D\in\mathcal D_M\), (1.3) and linearity of the finite sum give
\[
 \rho_M^{(0)}(D)=\lim_N\omega_N(\alpha_M^N(D)).                 \tag{2.4}
\]
Because \(\alpha_M^N\) is injective and therefore isometric,
\[
 |\rho_M^{(0)}(D)|
 \le\limsup_N\|\alpha_M^N(D)\|=\|D\|.                          \tag{2.5}
\]
If \(D\in\mathcal D_M\), then \(D^*D\in\mathcal D_M\), and
\[
 \rho_M^{(0)}(D^*D)
 =\lim_N\omega_N(\alpha_M^N(D)^*\alpha_M^N(D))\ge0.            \tag{2.6}
\]
Also \(\rho_M^{(0)}(I)=1\).

\LS{1}{3}{Extend to the coarse \(C^*\)-algebra.}

Equations (1.2) and (2.5) give a unique bounded extension
\(\rho_M:\W_M\to\C\).  Positivity persists under norm limits: if
\(A\in\W_M\) and \(D_j\to A\), then \(D_j^*D_j\to A^*A\), so
\[
 \rho_M(A^*A)=\lim_j\rho_M^{(0)}(D_j^*D_j)\ge0.                \tag{2.7}
\]
Together with \(\rho_M(I)=1\), this makes \(\rho_M\) a state.

\LS{1}{4}{Upgrade convergence from Weyl polynomials to every coarse observable.}

Fix \(A\in\W_M\) and \(\delta>0\).  By (1.2), choose
\(D\in\mathcal D_M\) with \(\|A-D\|<\delta\).  Then
\begin{align}
 |\omega_N(\alpha_M^N(A))-\rho_M(A)|
 &\le|\omega_N(\alpha_M^N(A-D))|\\
 &\quad+|\omega_N(\alpha_M^N(D))-\rho_M(D)|
       +|\rho_M(D-A)|\\
 &\le2\delta+
 |\omega_N(\alpha_M^N(D))-\rho_M(D)|.                          \tag{2.8}
\end{align}
The last term tends to zero by (2.4).  Therefore
\[
 \limsup_N|\omega_N(\alpha_M^N(A))-\rho_M(A)|\le2\delta.       \tag{2.9}
\]
Letting \(\delta\downarrow0\) proves (1.4).

\LS{1}{5}{Prove projective consistency.}

Fix \(M\le M'\) and \(A\in\W_M\).  For every \(N\ge M'\),
\[
 \alpha_{M'}^N(\alpha_M^{M'}(A))=\alpha_M^N(A).                \tag{2.10}
\]
Apply \(\omega_N\), take \(N\to\infty\), and use (1.4) at scales
\(M,M'\):
\[
 \rho_{M'}(\alpha_M^{M'}(A))=\rho_M(A).                        \tag{2.11}
\]
This is (1.5).

\LS{1}{6}{Construct the inductive-limit state.}

On the algebraic union prescribe
\[
 \Omega_0(\alpha_M^\infty(A)):=\rho_M(A).                      \tag{2.12}
\]
Equation (2.11) makes this well-defined.  It is positive, unital, and bounded by
one on each embedded algebra.  It therefore extends uniquely to the norm completion
\(\W_\infty\), proving (1.6). \QedStep

\section{Exact continuum-identification criterion}

\begin{corollary}[identification on one common algebra]\label{cor:identify}
Let \(\omega_{\rm ct}\) be a state on a continuum Weyl algebra \(\W_{\rm ct}\),
and let
\[
 \beta:\W_\infty\longrightarrow\W_{\rm ct}                    \tag{3.1}
\]
be an injective \( * \)-homomorphism satisfying
\(\beta\alpha_M^\infty=\beta_M\).  If
\[
 \chi_M(\xi)=\omega_{\rm ct}(\beta_M(W_M(\xi)))                \tag{3.2}
\]
for every \(M,\xi\), then
\[
 \Omega_\infty=\omega_{\rm ct}\circ\beta
 \quad\hbox{on }\W_\infty.                                    \tag{3.3}
\]
\end{corollary}

\begin{proof}
Equation (3.2) and Theorem \ref{thm:weylcriterion} give equality on every Weyl
generator, hence on \(\mathcal D_M\), hence by norm continuity on every
\(\W_M\).  The union of the embedded \(\W_M\) is norm dense in
\(\W_\infty\), so the equality extends to all of \(\W_\infty\).
\end{proof}

\begin{remark}[the exact remaining interacting statement]
For the finite-torus Wick-quartic states from the model sheet, the unresolved scalar
statement is
\[
 \lim_{N\to\infty}
 \omega_{N,L,\lambda,1}(\alpha_M^N(W_M(\xi)))
 =
 \omega_{P(\phi)_2,L}(\beta_M(W_M(\xi)))                       \tag{3.4}
\]
for every fixed \(M\) and \(\xi\).  Proving (3.4) closes the static,
full-sequence, identified fixed-coarse theorem by Corollary \ref{cor:identify}.
Neither Theorem \ref{thm:weylcriterion} nor Corollary \ref{cor:identify}
asserts (3.4).
\end{remark}

\section{Compatible symmetry}

\begin{proposition}
Suppose \(\vartheta_N\in\operatorname{Aut}(\W_N)\) satisfy
\[
 \vartheta_N\alpha_M^N=\alpha_M^N\vartheta_M,\qquad
 \omega_N\vartheta_N=\omega_N.                                \tag{4.1}
\]
Then every \(\rho_M\) and \(\Omega_\infty\) constructed above is invariant under
the induced symmetry.
\end{proposition}

\begin{proof}
For \(A\in\W_M\),
\[
 \rho_M(\vartheta_MA)
 =\lim_N\omega_N(\alpha_M^N(\vartheta_MA))
 =\lim_N\omega_N(\vartheta_N(\alpha_M^N A))
 =\rho_M(A).                                                    \tag{4.2}
\]
The inductive-limit assertion follows first on the algebraic union and then by norm
continuity.
\end{proof}

\section{Why the cited anharmonic LR theorem does not cover the Wick quartic}

Nachtergaele--Raz--Schlein--Sims, Theorem 4.1, assumes for the on-site
perturbation \(V\) that
\[
 V\in C^1(\R),\qquad V'\in L^1(\R),\qquad
 \kappa_V:=\int_{\R}|w|\,|\widehat{V'}(w)|\dd w<\infty.         \tag{5.1}
\]

\begin{proposition}[exact hypothesis failure]\label{prop:NRSSfail}
For every \(\lambda>0\) and \(c\ge0\), the Wick quartic
\[
 V_c(u):=\lambda(u^4-6cu^2+3c^2)                              \tag{5.2}
\]
does not satisfy \(V_c'\in L^1(\R)\).  Hence the cited theorem cannot be
applied to this interaction.
\end{proposition}

\subsection*{Proof in Lamport form}

\LS{1}{1}{Differentiate the stated polynomial.}

\[
 V_c'(u)=4\lambda u^3-12\lambda cu
 =4\lambda u(u^2-3c).                                         \tag{5.3}
\]

\LS{1}{2}{Obtain a lower bound on the tails.}

If \(|u|\ge\sqrt{6c}\), then
\[
 |u^2-3c|\ge\frac12u^2,                                       \tag{5.4}
\]
and therefore
\[
 |V_c'(u)|\ge2\lambda|u|^3.                                   \tag{5.5}
\]

\LS{1}{3}{Integrate.}

For every \(R\ge\sqrt{6c}\),
\[
 \int_R^\infty|V_c'(u)|\dd u
 \ge2\lambda\int_R^\infty u^3\dd u=\infty.                    \tag{5.6}
\]
Thus \(V_c'\notin L^1(\R)\), so the first analytic hypothesis in (5.1)
already fails. \QedStep

\begin{remark}
Proposition \ref{prop:NRSSfail} excludes only this citation route.  It is not a
nonexistence theorem for quartic Lieb--Robinson bounds.  A replacement would require a
polynomial-specific energy-weighted argument, with constants controlled uniformly in
the resolution.
\end{remark}

\section{Revised proof obligations}

\begin{longtable}{P{0.33\textwidth}P{0.17\textwidth}P{0.41\textwidth}}
\toprule
Item & Status & Exact content\\
\midrule
\endhead
Projective cofinal-subnet existence & \Pass & Proved independently in the compactness
companion.\\
Full-sequence reduction to Weyl generators & \Pass & Theorem
\ref{thm:weylcriterion}.\\
Continuum identification once generator limits are known & \Pass & Corollary
\ref{cor:identify}.\\
Interacting characteristic limit (3.4) & \Fail & This is the remaining static
finite-torus theorem.  It requires a uniform Euclidean or cluster-expansion comparison.\\
\(N\)-uniform gap and local second moment & \Fail & Not needed for the static weak-star
criterion, but still needed by the quantitative LPPL route.\\
Quartic interacting LR estimate & \Fail & The cited theorem assumes (5.1), which the
Wick quartic violates by Proposition \ref{prop:NRSSfail}.\\
Local-state-norm rate & \Fail & Pointwise convergence on a norm-dense set does not imply
uniform convergence over the norm-one ball.\\
\bottomrule
\end{longtable}

\begin{thebibliography}{9}
\bibitem{NRSS}
B.~Nachtergaele, H.~Raz, B.~Schlein, R.~Sims,
\emph{Lieb--Robinson bounds for harmonic and anharmonic lattice systems},
Commun. Math. Phys. \textbf{286} (2009), 1073--1098.
\end{thebibliography}

\end{document}
